\section{Conclusion}\label{sec:conclusion}

This study audited the operation implementing a human-grounded visual-quality
criterion in one multi-video summarization pipeline. The deployed
SIFT+BoVW+TF-IDF score factors into descriptor count and mean IDF weight. Under the
evaluated vocabularies, normalizations and extractors, descriptor count dominates:
the canonical score and keypoint density have statistically equivalent
MOS/DMOS associations within the operational $\pm.02$ margin on KADID and
KonIQ, whereas the deployed formulation produces a near-exact sign-reversed
count ordering in the large-database regime. In the broadcast domain,
the deployed score tracks its own descriptor count at $+.994$ and remains
bounded when SIFT is replaced by 19 convolutional extractors.

The mechanism predicts a consequential failure. Across 800 unmodified--degraded
broadcast pairs, the deployed operation selects the construction-defined
higher-fidelity member in 65.1\% of decisions and systematically prefers
white-noise and JPEG-degraded copies. LIQE reaches 99.5\%. A dependency-aware
screening rule retains LIQE, QualiCLIP and CLIP-IQA$+$; substituting the complete
visual-quality operation with a LIQE-based operation, while holding every other
pipeline stage fixed, changes 22 of 42 selected segments. The DOVER
technical-branch global-range point estimate increases by 25.5 percentage
points, but both topic-block and hierarchical intervals cross zero, so no
population-level output effect is confirmed.

The supported conclusion is therefore specific: the audited operation is not a
reliable perceptual-quality proxy in the tested regimes, and eligible NR-IQA
replacements materially alter its decisions. With frames and pooling matched,
SIFT+BoVW and LIQE still disagree in 16 of 42 clusters; 13 of these disagreements
overlap the 22 changes in the complete-operation comparison. The study does not establish that
all BoVW formulations fail, that redundancy clusters contain semantically
identical alternatives, or that viewers prefer the modified summaries. Human
pairwise evaluation and broader video-native quality assessment remain the next
tests needed to connect technical visual quality to compositional value.
